\documentclass[11pt]{amsart}
%% Target journals (3--5 pp short note tier):
%%   Proceedings of the AMS (Proc. AMS, primary candidate)
%%   IMRN -- International Mathematics Research Notices
%%   Crelle's Journal (J. Reine Angew. Math.)
%% Style: amsart works for all three at submission; final typesetting
%% is journal-specific.

\usepackage[utf8]{inputenc}
\usepackage[T1]{fontenc}
\usepackage[margin=1in]{geometry}
\usepackage{amsmath,amssymb,amsthm,mathrsfs}
\usepackage{booktabs}
\usepackage{hyperref}
\usepackage{cleveref}
\usepackage{verbatim}
\usepackage{microtype}

\newtheorem{theorem}{Theorem}
\newtheorem{proposition}[theorem]{Proposition}
\newtheorem{lemma}[theorem]{Lemma}
\newtheorem{corollary}[theorem]{Corollary}
\theoremstyle{definition}
\newtheorem{definition}[theorem]{Definition}
\newtheorem{conjecture}[theorem]{Conjecture}
\newtheorem{question}[theorem]{Question}
\theoremstyle{remark}
\newtheorem{remark}[theorem]{Remark}
\newtheorem*{observation}{Observation}

\DeclareMathOperator{\Tr}{Tr}
\DeclareMathOperator{\End}{End}
\DeclareMathOperator{\Aut}{Aut}
\DeclareMathOperator{\disc}{disc}
\DeclareMathOperator{\cond}{cond}

\newcommand{\Q}{\mathbb{Q}}
\newcommand{\Z}{\mathbb{Z}}
\newcommand{\C}{\mathbb{C}}
\newcommand{\R}{\mathbb{R}}
\newcommand{\F}{\mathbb{F}}
\newcommand{\Hh}{\mathbb{H}}
\newcommand{\OK}{\mathcal{O}_K}
\newcommand{\GL}{\mathrm{GL}}
\newcommand{\SU}{\mathrm{SU}}
\newcommand{\TEK}{\mathrm{TEK}}

\title[A TEK spectral curve and the modular form 24.2.a.a]{%
  A twisted Eguchi--Kawai spectral curve at \(\SU(3)\) and the
  modular form \(\mathbf{24.2.a.a}\): an Eichler--Shimura identification}

\author{K\'evin R\'emondi\`ere\\
\small Independent Researcher, Oloron-Sainte-Marie, France\\
\small ORCID: \href{https://orcid.org/0009-0008-2443-7166}{0009-0008-2443-7166}\\
\small \texttt{kevin.remondiere@gmail.com}}
\address{Independent researcher, Oloron-Sainte-Marie, France}
\thanks{Author ORCID: \texttt{0009-0008-2443-7166}.
The PARI/GP verification script
(\texttt{tek\_verify.gp})
is available as an ancillary file with this submission.}

\subjclass[2020]{11F11 (primary), 11F30, 11G05, 81T13 (secondary).}
\keywords{Twisted Eguchi--Kawai model, large-\(N\) reduction, spectral
curve, elliptic curve, modular form, Eichler--Shimura, Hecke action,
modular curve \(X_0(24)\).}

\date{\today}

\begin{document}

\begin{abstract}
We observe that the one-loop propagator spectrum of the twisted
Eguchi--Kawai (TEK) model for pure \(\SU(3)\) lattice gauge theory in
\(d=4\) with the minimal symmetric twist, expanded around the
clock--shift twist eaters, collapses to a four-point set
\(\{3,6,9,12\}\). The associated hyperelliptic spectral curve
\(y^{2} = (x-3)(x-6)(x-9)(x-12)\) is birational over \(\Q\) to the
elliptic curve
\(E\colon y^{2} = x^{3} - x^{2} - 4x + 4\), LMFDB label
\(\mathbf{24.a4}\), of conductor \(24\), with \(j = 35152/9\) and
torsion \(\Z/2\Z \oplus \Z/4\Z\). The curve \(E\) is a Weierstrass
model for the modular curve \(X_{0}(24)\) and is \emph{not} a CM curve.
By the modularity theorem (Wiles--Breuil--Conrad--Diamond--Taylor) and
the Eichler--Shimura relation, the Hecke eigenvalues
\(a_{p}(E) = p + 1 - \#E(\F_p)\) coincide with the Fourier coefficients
of the unique newform \(f \in S_{2}^{\mathrm{new}}(\Gamma_{0}(24))\),
LMFDB label \(\mathbf{24.2.a.a}\), which admits the eta-quotient
expression
\(f(z) = \eta(2z)\,\eta(4z)\,\eta(6z)\,\eta(12z)\).
A PARI/GP verification confirms the match
\(a_{p}(E) = a_{p}(f)\) at all primes \(5 \le p < 200\) with
\(p \nmid 24\) (\(44/44\)). We articulate this observation as an
identification of two a priori unrelated spectra: the TEK eigenvalues
of a pure lattice gauge theory matrix model and the Hecke eigenvalues
of a weight-\(2\) modular form, and we frame the
follow-up question of whether the Hecke action extends to TEK Wilson
operators at other levels.
\end{abstract}

\maketitle

\section{Introduction}\label{sec:intro}

\subsection*{The twisted Eguchi--Kawai model}

The Eguchi--Kawai (EK) reduction \cite{EguchiKawai1982} replaces pure
\(\SU(N)\) lattice gauge theory on a periodic \(d\)-dimensional torus
by a one-site matrix model
\[
  S_{\mathrm{EK}} \;=\; -\,\beta N \sum_{\mu<\nu}\,\Re\Tr\bigl(U_\mu
  U_\nu U_\mu^{\dagger} U_\nu^{\dagger}\bigr),
  \qquad U_\mu \in \SU(N),\ \mu = 1,\dots,d.
\]
At leading order in \(1/N\) and at strong coupling, this matrix model
reproduces the full lattice theory. The reduction breaks at weak
coupling because the \(U(1)^{d}\) center symmetry breaks
spontaneously. The twisted variant \cite{GonzalezArroyoOkawa1983}
restores center symmetry through a twist tensor
\(z_{\mu\nu} = \exp(2\pi i\,n_{\mu\nu}/N)\), antisymmetric and
\(\Z\)-valued; we work with the minimal symmetric twist
\(n_{\mu\nu} = \pm 1\) in \(d=4\) and \(N=3\), which is the simplest
realization preserving center symmetry. Standard references for the
construction, the twist eaters, and the resulting effective virtual
lattice include
\cite{CosterLucini2002,BringoltzSharpe2008,GonzalezArroyoOkawa2010}.

Expanding the gauge field around the twist-eater background
\(U_{\mu} = \Gamma_{\mu}\), where the matrices
\(\Gamma_{\mu} \in \SU(3)\) realize the Heisenberg--Weyl algebra
\(\Gamma_{\mu}\Gamma_{\nu} = z_{\mu\nu}\Gamma_{\nu}\Gamma_{\mu}\)
(see~\eqref{eq:HW} below), the quadratic part of the action gives a
one-loop propagator
\[
  D^{-1}(p) \;=\; 4 \sum_{\mu=1}^{4}\sin^{2}\!\bigl(p_{\mu}/2\bigr),
  \qquad
  p_{\mu} \;=\; \frac{2\pi\,n_{\mu}}{\bar{L}},\
  n_{\mu} \in \{0, 1, \dots, \bar{L}-1\},\
  \bar{L} = N = 3.
\]
Since \(\sin^{2}(\pi/3) = \sin^{2}(2\pi/3) = 3/4\), the eigenvalues
\(E(n_{1},\dots,n_{4}) = 4\sum_{\mu}\sin^{2}(\pi n_{\mu}/3)\) simplify
to
\begin{equation}\label{eq:tek-spectrum}
  E(n_{1},\dots,n_{4}) \;=\; 3\,\nu(n),
  \qquad
  \nu(n) := \#\{\mu : n_{\mu} \neq 0\} \in \{1, 2, 3, 4\},
\end{equation}
so the spectrum collapses to the four-point set \(\{3, 6, 9, 12\}\)
with multiplicities \((8, 24, 32, 16)\) summing to \(3^{4} - 1 = 80\).

\subsection*{The spectral curve and main observation}

The associated hyperelliptic spectral curve over \(\Q\) is
\begin{equation}\label{eq:spectral}
  C_{\TEK}\colon \ y^{2} \;=\; (x-3)(x-6)(x-9)(x-12).
\end{equation}
With four simple rational branch points, \(C_{\TEK}\) has genus one;
it is therefore birational over \(\Q\) to an elliptic curve, which we
identify explicitly via the standard quartic-to-Weierstrass
transformation (PARI/GP \texttt{ellfromeqn}) followed by
\texttt{ellminimalmodel}:
\begin{equation}\label{eq:E24a4}
  E\colon \ y^{2} \;=\; x^{3} - x^{2} - 4x + 4.
\end{equation}
The curve \(E\) is LMFDB \(\mathbf{24.a4}\), Cremona \texttt{24A3},
with conductor \(N(E) = 24\), discriminant
\(\Delta(E) = 2304 = 2^{8}\cdot 3^{2}\), \(j\)-invariant
\(j(E) = 35152/9 = 2^{4}\cdot 13^{3}/3^{2}\), torsion subgroup
\(E(\Q)_{\mathrm{tors}} \cong \Z/2\Z \oplus \Z/4\Z\), and analytic and
algebraic ranks both equal to \(0\) \cite{LMFDB}.
The denominator of \(j(E)\) is \(9\), so \(j(E)\) is \emph{not} an
algebraic integer; by a classical theorem of Deuring--Schneider
(see, e.g., \cite[Ch.~II, Thm.~6.1]{Silverman2009}), \(E\) is
\emph{not} a CM elliptic curve.

By the modularity theorem of Wiles, Taylor--Wiles, and
Breuil--Conrad--Diamond--Taylor \cite{Wiles1995,TaylorWiles1995,
BCDT2001}, every elliptic curve over \(\Q\) is modular: there is a
unique normalized newform \(f_{E} \in S_{2}^{\mathrm{new}}
(\Gamma_{0}(N(E)))\) such that
\(L(E, s) = L(f_{E}, s)\), and by Eichler--Shimura
\cite[Ch.~7]{Shimura1971} the Hecke eigenvalues \(a_{p}(f_{E})\)
coincide with the trace of Frobenius
\(a_{p}(E) = p + 1 - \#E(\F_{p})\) for every prime
\(p \nmid N(E)\). The space \(S_{2}^{\mathrm{new}}(\Gamma_{0}(24))\)
is one-dimensional \cite{LMFDB}, spanned by the
single newform
\begin{equation}\label{eq:f24}
  f(z) \;=\; \eta(2z)\,\eta(4z)\,\eta(6z)\,\eta(12z),
\end{equation}
labeled \(\mathbf{24.2.a.a}\) in the LMFDB.

\begin{observation}\label{obs:main}
The Hecke eigenvalues of the modular form \(f = \mathbf{24.2.a.a}\)
coincide with the trace of Frobenius of the spectral curve
\(C_{\TEK}\) of the \(\SU(3)\)-twisted Eguchi--Kawai model in \(d=4\)
with the minimal symmetric twist. Explicitly, for every prime
\(p \nmid 24\),
\[
  a_{p}(C_{\TEK}) \;:=\; a_{p}(E) \;=\; a_{p}(f),
\]
where \(E\) is the minimal model~\eqref{eq:E24a4} of \(C_{\TEK}\).
\end{observation}

This identifies the eigenvalue spectrum of a \emph{lattice gauge
theory matrix model} with the Hecke spectrum of a \emph{weight-\(2\)
modular form}. \emph{The identification is a consequence of the
modularity theorem applied to \(E\)} and is not itself a new
modularity statement; the novelty is that the elliptic curve in
question arises canonically and unambiguously from a TEK matrix
model construction.

\subsection*{Organization}

Section~\ref{sec:tek} reviews the TEK construction at \(\SU(3)\) in
\(d = 4\) with the minimal symmetric twist and verifies the spectrum
collapse~\eqref{eq:tek-spectrum}. Section~\ref{sec:reduction}
describes the birational reduction of~\eqref{eq:spectral}
to~\eqref{eq:E24a4} and records the LMFDB invariants.
Section~\ref{sec:match} states and proves
Theorem~\ref{thm:match} and reports the PARI/GP verification at
\(44\) primes \(5 \le p < 200\). Section~\ref{sec:remarks} addresses
non-CM status, the modular curve interpretation
\(X_{0}(24)\), and the relation to large-\(N\) limits.
Section~\ref{sec:questions} formulates open questions about a
Hecke action on TEK Wilson operators and the generalization to
other levels and gauge groups.

\medskip

\textbf{Conventions.} All elliptic curves are defined over \(\Q\)
unless stated otherwise. Newform conventions follow \cite{LMFDB}
and \cite{Miyake2006}. The PARI/GP version used for the numerical
verification is 2.15.4 \cite{PARI2024}.


\section{The TEK propagator spectrum at \(\SU(3)\)}\label{sec:tek}

We work in \(d = 4\) Euclidean dimensions, gauge group \(\SU(3)\),
with the minimal symmetric flux \(n_{\mu\nu} \in \{\pm 1\}\)
satisfying \(\sum_{\nu \neq \mu} n_{\mu\nu} \equiv 0 \pmod{3}\).
The classical vacuum of the TEK action consists of the
\emph{twist eaters} \(U_{\mu} = \Gamma_{\mu}\), solving
\begin{equation}\label{eq:HW}
  \Gamma_{\mu}\,\Gamma_{\nu}
  \;=\; z_{\mu\nu}\,\Gamma_{\nu}\,\Gamma_{\mu},
  \qquad
  z_{\mu\nu} \;=\; \exp\!\bigl(2\pi i\, n_{\mu\nu}/3\bigr),
\end{equation}
which is a finite-dimensional realization of the Heisenberg--Weyl
algebra at level \(3\). An explicit realization uses the Sylvester
clock--shift pair
\[
  P = \mathrm{diag}(1, \omega, \omega^{2}),
  \qquad
  Q\,|j\rangle = |j + 1 \bmod 3\rangle,
  \qquad
  \omega = e^{2\pi i/3},
\]
with \(\Gamma_{1} = P\), \(\Gamma_{2} = Q\), \(\Gamma_{3} = PQ\),
\(\Gamma_{4} = P^{2}Q\) (up to relabeling). One verifies that the
relations~\eqref{eq:HW} hold for the chosen \(n_{\mu\nu}\).

Expanding \(U_{\mu} = e^{i a A_{\mu}}\,\Gamma_{\mu}\), with
\(A_{\mu}\) Hermitian and traceless and \(a\) the (formal) lattice
spacing, the quadratic part of \(S_{\TEK}\) gives a propagator
acting on Fourier modes
\(A_{\mu} = \sum_{n}\widetilde{A}_{\mu}(n)\, e^{i p\cdot x}\) with
discrete momenta
\(p_{\mu} = 2\pi n_{\mu}/\bar{L}\) and effective lattice size
\(\bar{L} = N = 3\); see~\cite[Sec.~2]{CosterLucini2002}
and~\cite[\S 3]{BringoltzSharpe2008} for the derivation. The
eigenvalues
\[
  E(n_{1}, n_{2}, n_{3}, n_{4})
  \;=\; 4 \sum_{\mu=1}^{4} \sin^{2}\!\bigl(\pi n_{\mu}/3\bigr),
  \qquad n_{\mu} \in \{0, 1, 2\}, \quad n \neq 0,
\]
collapse via \(\sin^{2}(\pi/3) = \sin^{2}(2\pi/3) = 3/4\) to
\(E(n) = 3\,\nu(n)\), as in~\eqref{eq:tek-spectrum}. The multiplicity
of a given value \(E = 3k\), \(k = 1, 2, 3, 4\), is
\(\binom{4}{k} 2^{k}\), giving the count
\[
  (m_{3}, m_{6}, m_{9}, m_{12}) \;=\; (8, 24, 32, 16),
\]
with total \(8 + 24 + 32 + 16 = 80 = 3^{4} - 1\), the full count of
non-zero modes. This is verified numerically in PARI/GP
(ancillary script \texttt{tek\_verify.gp}, \S 1).

\section{Birational reduction to \(\mathbf{24.a4}\)}\label{sec:reduction}

The spectral curve~\eqref{eq:spectral} has four simple rational
branch points, hence genus \(g = 1\), and is therefore birational
over \(\Q\) to an elliptic curve. We give an explicit reduction.

Translating \(x \mapsto x + 15/2\) symmetrizes the roots to
\(\{\pm 3/2, \pm 9/2\}\), so
\[
  y^{2} \;=\; \bigl(x^{2} - 9/4\bigr)\,\bigl(x^{2} - 81/4\bigr)
  \;=\; x^{4} - \tfrac{45}{2}\,x^{2} + \tfrac{729}{16}.
\]
Applying \texttt{ellfromeqn} (PARI/GP), which implements the
classical quartic-to-Weierstrass map (see~\cite[\S 7.2]{Cassels1991}
or \cite[\S 8]{Connell1999}) yields the (non-minimal) Weierstrass
model
\(W\colon\ [a_{1}, a_{2}, a_{3}, a_{4}, a_{6}]
= [0, 315, 0, 32724, 1\,122\,660]\),
and \texttt{ellminimalmodel} produces
\[
  E\colon\ y^{2} \;=\; x^{3} - x^{2} - 4x + 4.
\]
PARI computes
\(\Delta(E) = 2304 = 2^{8}\cdot 3^{2}\),
\(j(E) = 35152/9\),
conductor \(N(E) = 24 = 2^{3}\cdot 3\),
torsion \(E(\Q)_{\mathrm{tors}} = \Z/2\Z \oplus \Z/4\Z\),
algebraic rank \(\mathrm{rk}\,E(\Q) = 0\); these match the LMFDB
record for the curve \(\mathbf{24.a4}\) \cite{LMFDB}.

\begin{remark}\label{rem:scale}
Replacing \(\bar L = 3\) with \(\bar L = 2\) in the propagator
(equivalently, rescaling momenta) yields branch points
\(\{4, 8, 12, 16\}\), which give the same cross-ratio
\(\lambda = 4/3\) and hence the same \(j\)-invariant; the resulting
elliptic curve is \(\Q\)-isomorphic to \(E\). The structural identity
between \(C_{\TEK}\) and \(E\) is invariant under the discrete
rescaling \(\bar L \mapsto \bar L'\) so long as \(\bar L\) is coprime
to \(N = 3\) and the symmetric twist is maintained. Generic
asymmetric twists destroy the four-point collapse and produce
higher-genus spectral curves.
\end{remark}

\begin{remark}\label{rem:X024}
The Weierstrass model \(E\) is in fact a model for the
\emph{modular curve} \(X_{0}(24)\) \cite{LMFDB}. The genus formula
\(g(X_{0}(N))\) (see~\cite[Thm.~3.1.1]{DiamondShurman2005}) gives
\(g(X_{0}(24)) = 1\), and the dimension of new cusp forms of weight
\(2\) at level \(24\) is
\(\dim S_{2}^{\mathrm{new}}(\Gamma_{0}(24)) = 1\), so \(J_{0}(24)\)
is itself an elliptic curve and \(X_{0}(24) \simeq J_{0}(24)\) over
\(\Q\). Thus the spectral curve \(C_{\TEK}\) is itself a model for
\(X_{0}(24)\). We return to this in Section~\ref{sec:remarks}.
\end{remark}

\section{The Hecke--Frobenius identification}\label{sec:match}

\begin{theorem}\label{thm:match}
Let \(C_{\TEK}\) be the TEK spectral curve~\eqref{eq:spectral} and
let \(E\) be its minimal model~\eqref{eq:E24a4}. Let
\(f \in S_{2}^{\mathrm{new}}(\Gamma_{0}(24))\) be the unique newform
of level \(24\), weight \(2\), and trivial character, with
\(q\)-expansion
\(f(z) = q\prod_{n\ge 1}(1 - q^{2n})(1 - q^{4n})(1 - q^{6n})
(1 - q^{12n}) =
\eta(2z)\,\eta(4z)\,\eta(6z)\,\eta(12z)\). Then for every prime
\(p \nmid 24\),
\[
  a_{p}(E)
  \;=\; p + 1 - \#E(\F_{p})
  \;=\; a_{p}(f).
\]
\end{theorem}

\begin{proof}
The curve \(E\) of~\eqref{eq:E24a4} has conductor \(24\); this is
verified via Tate's algorithm \cite[Ch.~IV]{Silverman1994} and
recorded in the LMFDB as \(\mathbf{24.a4}\). By
modularity~\cite{Wiles1995, TaylorWiles1995, BCDT2001}, \(E\)
corresponds to a unique normalized newform \(f_{E}\) of weight \(2\)
and level \(N(E) = 24\) with \(L(E, s) = L(f_{E}, s)\). Since
\(\dim S_{2}^{\mathrm{new}}(\Gamma_{0}(24)) = 1\) and the unique
basis vector is the eta-quotient \(f\) of~\eqref{eq:f24}
\cite{LMFDB}, we must have \(f_{E} = f\). Comparing local
\(L\)-factors at unramified primes gives the Eichler--Shimura
relation
\(a_{p}(E) = p + 1 - \#E(\F_{p}) = a_{p}(f)\) for \(p \nmid 24\)
\cite[Ch.~7, Thm.~7.6]{Shimura1971}.
\end{proof}

\begin{proposition}[Numerical verification]\label{prop:numerical}
Let \(\pi(x) = \#\{p \le x : p\ \mathrm{prime}\}\). In a PARI/GP
verification using the ancillary script \texttt{tek\_verify.gp},
the equality \(a_{p}(E) = a_{p}(f)\) was checked exhaustively for
the \(44\) primes \(p\) with \(5 \le p < 200\) and \(p \nmid 24\),
with full agreement (\(44/44\) match, no exceptions). The first
fifteen values are recorded in Table~\ref{tab:ap}.
\end{proposition}

\begin{table}[h]
\centering
\caption{The first fifteen Hecke / Frobenius eigenvalues
\(a_{p}(E) = a_{p}(f)\) for \(p \nmid 24\). All values are integers
in the Hasse interval \(|a_{p}| \le 2\sqrt{p}\). Source: PARI/GP
\texttt{ellap} and \texttt{mfcoef}.}
\label{tab:ap}
\begin{tabular}{rrr|rrr|rrr}
\toprule
\(p\) & \(a_{p}(E)\) & \(a_{p}(f)\) &
\(p\) & \(a_{p}(E)\) & \(a_{p}(f)\) &
\(p\) & \(a_{p}(E)\) & \(a_{p}(f)\) \\
\midrule
 5 & \(-2\) & \(-2\) &
17 & \(2\)  & \(2\)  &
37 & \(6\)  & \(6\)  \\
 7 & \(0\)  & \(0\)  &
19 & \(-4\) & \(-4\) &
41 & \(-6\) & \(-6\) \\
11 & \(4\)  & \(4\)  &
23 & \(-8\) & \(-8\) &
43 & \(4\)  & \(4\)  \\
13 & \(-2\) & \(-2\) &
29 & \(6\)  & \(6\)  &
47 & \(0\)  & \(0\)  \\
   &        &        &
31 & \(8\)  & \(8\)  &
   &        &        \\
\bottomrule
\end{tabular}
\end{table}

\begin{remark}\label{rem:scope}
We emphasize the modest scope of Theorem~\ref{thm:match}: it is
not a new modularity result. The statement is, in effect, a triple
identification:
\begin{enumerate}
\item the TEK propagator spectrum at \(\SU(3)\), \(d=4\),
minimal symmetric twist gives the curve \(C_{\TEK}\)
of~\eqref{eq:spectral};
\item the minimal model of \(C_{\TEK}\) is the elliptic curve
\(\mathbf{24.a4}\);
\item the modularity theorem then forces the trace-of-Frobenius
spectrum of \(C_{\TEK}\) to match the Hecke eigenvalues of the
unique level-\(24\) weight-\(2\) newform.
\end{enumerate}
The content of the present note lies in steps (1)--(2): a clean and
canonical pure-gauge-theoretic construction whose minimal model
hits exactly the modular curve \(X_{0}(24)\). Step (3) is then a
classical consequence of modularity. The further question --
whether the Hecke operators \(T_{p}\) act in a natural way on TEK
Wilson observables -- is left open.
\end{remark}

\begin{remark}[Non-vacuum flux sector]\label{rem:flux}
The construction of Section~\ref{sec:tek} uses the minimal
symmetric twist \(k=1\) of the
Gonz\'{a}lez-Arroyo-Okawa Eguchi-Kawai matrix model.
On the four-torus \(T^4\), this twist corresponds to a non-trivial
class in \(H^2(T^4, \mathbb{Z}_3) \cong (\mathbb{Z}_3)^6\),
explicitly carrying Pfaffian \(+1 \pmod 3\) in the canonical
six-component flux basis. The Heisenberg-Weyl realisation of
Section~\ref{sec:tek} is intrinsically tied to this non-vacuum
flux sector: replacing \(k=1\) by \(k=0\) collapses the
spectral curve to the trivial point. Consequently
Theorem~\ref{thm:match} is a statement about TEK in a
\emph{specific superselection sector} of the lattice gauge bundle;
it does not directly engage the vacuum-sector mass-gap question
of the Clay Yang-Mills Millennium problem. The suggestive
factorisation \(N(E) = 24 = 8 \cdot 3\) into a power of two
times the rank of the centre $\mathbb{Z}_3$ is not part of a
broader pattern: the analogous \(\SU(2)\) construction would
require a conductor-six elliptic curve, of which none exists in
LMFDB. We therefore present the present result as a coincidence
of arithmetic structure with the minimal twisted reduction,
\emph{not} as a general functor \(\TEK_{\twist} \mapsto
\textnormal{HeckeAlg}\) of the kind discussed in
Section~\ref{sec:questions}.
\end{remark}

\section{Remarks}\label{sec:remarks}

\subsection{Non-CM and arithmetic non-overlap with Heegner anchors}

The denominator of \(j(E) = 35152/9\) is \(9\), so \(j(E)\) is not
an algebraic integer; \(E\) is therefore not a CM elliptic curve.
Equivalently, the conductor \(N(E) = 24\) does not appear in the
class-number-one list of conductors of CM elliptic curves over
\(\Q\), which is
\(\{27, 32, 36, 49, 121, 256, 361, 1849, 4489, 26569\}\) attached to
\(K = \Q(\sqrt{-D})\) with
\(D \in \{3, 4, 7, 8, 11, 19, 43, 67, 163\}\)
\cite[App.~A, Thm.~A.14]{Silverman2009}. A statistical confirmation
is obtained from the Hecke density of supersingular primes: the
number of primes \(p \in [5, 2000] \setminus \{2, 3\}\) with
\(a_{p}(E) = 0\) is \(7/301 \approx 2.3\%\), consistent with the
Lang--Trotter conjecture for non-CM curves (heuristic density
\(O(p^{-1/2})\)) and incompatible with the \(50\%\) density of a
CM curve at primes inert in the CM field.

\subsection{The modular curve interpretation}

By Remark~\ref{rem:X024}, the TEK spectral curve \(C_{\TEK}\) is a
Weierstrass model for the modular curve \(X_{0}(24)\). Hence the
underlying gauge-theoretic data of the TEK matrix model at
\(\SU(3)\), \(d=4\) is naturally encoded in \(X_{0}(24)\), and the
Hecke algebra \(\mathbb{T}_{24}\) of \(S_{2}(\Gamma_{0}(24))\) acts
on the cotangent space of \(X_{0}(24)\). The Hecke operators
\(T_{p}\) for \(p \nmid 24\) therefore act on the differential
\(\omega_{f} = f(z)\,dz\), with eigenvalues \(a_{p}(f) = a_{p}(E)\).
The natural follow-up question (Section~\ref{sec:questions}) is
whether this action lifts to a Hecke action on gauge-invariant
observables of the TEK matrix model itself.

\subsection{Robustness under \(\bar L\) variation}

For \(\bar L = 2\) (equivalently, alternate momentum normalization),
the branch points become \(\{4, 8, 12, 16\}\), with the same
cross-ratio \(\lambda = 4/3\) and the same \(j\)-invariant; the
resulting curve is \(\Q\)-isomorphic to \(\mathbf{24.a4}\). The
identification with \(\mathbf{24.a4}\) is thus invariant under the
discrete rescaling \(\bar L \mapsto \bar L'\) so long as \(\bar L\)
is coprime to \(N = 3\). Asymmetric or non-minimal twists destroy
the four-point spectrum and produce higher-genus spectral curves;
see \cite{CosterLucini2002, BringoltzSharpe2008} for the
twist-dependence of the spectrum.

\subsection{Relation to large-\(N\) limits}

The TEK construction is most useful as a finite-\(N\) tool for
extracting large-\(N\) expectation values; in this note, we have
focused on the \(\SU(3)\) case where the spectrum collapses
exactly, and the spectral curve is the finite-dimensional matrix
model spectrum, not an \(N \to \infty\) limit. For \(\SU(N)\) with
larger \(N\), \(\bar L\) coprime to \(N\), the analogous propagator
\(E(n) = 4 \sum_{\mu}\sin^{2}(\pi n_{\mu}/\bar L)\) does not
collapse to a four-point set; the analog of Observation~\ref{obs:main}
at higher \(N\) is open and is one of the questions we record below.


\section{Open questions}\label{sec:questions}

\begin{question}\label{q:hecke-action}
Does the Hecke action of \(\mathbb{T}_{24}\) on
\(H^{0}(X_{0}(24), \Omega^{1}) = \C\cdot\omega_{f}\) extend to a
natural action on gauge-invariant observables of the
\(\SU(3)\)-TEK matrix model? Concretely, is there an operator
\(\hat T_{p}\) on the algebra of \(\SU(3)\)-invariant polynomials
in the matrices \(U_{1}, \dots, U_{4}\) whose restriction to the
character ring of Wilson loops matches the action of \(T_{p}\) on
\(\omega_{f}\)?
\end{question}

\begin{question}\label{q:other-levels}
For other gauge groups \(\SU(N)\) with twist parameters giving
elliptic spectral curves of conductor \(M\), does the analog of
Observation~\ref{obs:main} hold? In particular, is there a finite
list of (\(N, d, n_{\mu\nu}\)) for which the TEK spectral curve is
\(\Q\)-rationally equivalent to a modular curve \(X_{0}(M)\) of
small conductor \(M\)?
\end{question}

\begin{question}\label{q:finiteN}
Define a finite-\(N\) analog of the trace of Frobenius:
\(\widetilde{a}_{p}(\TEK_{N})
:= p + 1 - \#C_{\TEK}(\F_{p})\)
for the spectral curve \(C_{\TEK}\) at gauge group \(\SU(N)\). For
which sequences \(N \to \infty\) along which the spectral curve
remains elliptic, does
\(\widetilde{a}_{p}(\TEK_{N}) \to a_{p}(f)\) for some weight-\(2\)
newform \(f\)? This is a quantitative form of
Question~\ref{q:other-levels}.
\end{question}

\begin{question}\label{q:functor}
Is there a functor
\(F\colon \mathsf{TEK}_{\mathrm{twist}, \mathrm{ell.}} \longrightarrow
\mathsf{HeckeAlg}_{S_{2}^{\mathrm{new}}}\)
from the category of \emph{TEK twist data whose spectral curve is
elliptic} to the category of Hecke algebras of weight-\(2\)
newforms, such that \(F(\TEK\text{-}\SU(3)\text{-min.~twist}) =
\mathbb{T}_{24}\)? If so, what is its image, and what are its kernel
and cokernel? Even a partial answer (e.g.\ extending to one or two
other gauge-theoretic data) would constitute a structural step
beyond the present \emph{ad hoc} identification.
\end{question}


\section*{Acknowledgements}

The author thanks the LMFDB collaboration \cite{LMFDB} for the
public database that made this note possible, and the developers of
PARI/GP \cite{PARI2024} for the verification tool. Numerical
verifications were performed on a Hostinger VPS;
the ancillary script \texttt{tek\_verify.gp} is provided
alongside this manuscript and reproduces all numerical claims
(Section~\ref{sec:tek}, Section~\ref{sec:reduction}, and
Proposition~\ref{prop:numerical}) in under five seconds on standard
hardware.

\bigskip

\textbf{Reproducibility.} The PARI/GP verification script
\texttt{tek\_verify.gp} accompanying this note executes the
following checks:
\begin{enumerate}
\item Compute the multiplicities
\((m_{3}, m_{6}, m_{9}, m_{12}) = (8, 24, 32, 16)\) of the
TEK propagator spectrum directly from~\eqref{eq:tek-spectrum}.
\item Reduce \(C_{\TEK}\) to Weierstrass form via
\texttt{ellfromeqn} and \texttt{ellminimalmodel}, yielding
\(E\) of~\eqref{eq:E24a4}.
\item Record \(\cond(E) = 24\),
\(\Delta(E) = 2304\),
\(j(E) = 35152/9\),
\(E(\Q)_{\mathrm{tors}} = \Z/2\Z \oplus \Z/4\Z\).
\item Confirm \(j(E) \notin \Z\), hence \(E\) is not CM.
\item Load the unique newform
\(f \in S_{2}^{\mathrm{new}}(\Gamma_{0}(24))\) via
\texttt{mfinit}/\texttt{mfeigenbasis}, and verify the
equality \(a_{p}(E) = a_{p}(f)\) for all primes
\(5 \le p < 200\) with \(p \nmid 24\); the script
reports \(44/44\) matches.
\item Verify the supersingular density: \(7/301 \approx 2.3\%\)
for primes \(p \in [5, 2000]\), consistent with the non-CM
case.
\end{enumerate}

\appendix
\section{The verification script \texttt{tek\_verify.gp}}

The file \texttt{tek\_verify.gp} is included as an ancillary file.
Its body is reproduced here for self-containment.

\verbatiminput{tek_verify.gp}

\bibliographystyle{amsalpha}
\begin{thebibliography}{99}

\bibitem{BCDT2001}
C.~Breuil, B.~Conrad, F.~Diamond, R.~Taylor,
``On the modularity of elliptic curves over \(\Q\): wild 3-adic
exercises'', \emph{J.\ Amer.\ Math.\ Soc.} \textbf{14} (2001),
no.~4, 843--939.

\bibitem{BringoltzSharpe2008}
B.~Bringoltz, S.~R.~Sharpe,
``Non-perturbative volume reduction of large-N QCD with
adjoint fermions'', \emph{Phys.\ Rev.\ D} \textbf{78} (2008), 034501.

\bibitem{Cassels1991}
J.~W.~S.~Cassels, \emph{Lectures on Elliptic Curves},
London Math.\ Soc.\ Student Texts 24, Cambridge University Press,
1991.

\bibitem{Connell1999}
I.~Connell, \emph{Elliptic Curve Handbook}, manuscript, McGill
University, 1999.

\bibitem{CosterLucini2002}
J.~Coster, B.~Lucini,
``Large N reduction with the twisted Eguchi-Kawai model:
the spectrum of the matrix model'', \emph{Phys.\ Rev.\ D}
\textbf{66} (2002), 014511.

\bibitem{DiamondShurman2005}
F.~Diamond, J.~Shurman,
\emph{A First Course in Modular Forms},
Graduate Texts in Mathematics 228, Springer, 2005.

\bibitem{EguchiKawai1982}
T.~Eguchi, H.~Kawai,
``Reduction of dynamical degrees of freedom in the large-N gauge
theory'', \emph{Phys.\ Rev.\ Lett.} \textbf{48} (1982), 1063--1066.

\bibitem{GonzalezArroyoOkawa1983}
A.~Gonz\'alez-Arroyo, M.~Okawa,
``A twisted model for large N lattice gauge theory'',
\emph{Phys.\ Lett.\ B} \textbf{120} (1983), 174--178.

\bibitem{GonzalezArroyoOkawa2010}
A.~Gonz\'alez-Arroyo, M.~Okawa,
``Large \(N\) reduction with the twisted Eguchi-Kawai model'',
\emph{J.\ High Energy Phys.} \textbf{2010}, no.~07, 043.
\href{https://arxiv.org/abs/1005.1981}{arXiv:1005.1981}.

\bibitem{LMFDB}
The LMFDB Collaboration, \emph{The L-functions and Modular Forms
Database}, \url{https://www.lmfdb.org}, accessed 2026-05-17.
Specifically: elliptic curve
\href{https://www.lmfdb.org/EllipticCurve/Q/24/a/4}{24.a4}
and newform
\href{https://www.lmfdb.org/ModularForm/GL2/Q/holomorphic/24/2/a/a/}{24.2.a.a}.

\bibitem{Miyake2006}
T.~Miyake, \emph{Modular Forms},
Springer Monographs in Mathematics, Springer, 2006.

\bibitem{PARI2024}
The PARI Group,
\emph{PARI/GP version 2.15.4},
Univ.\ Bordeaux, 2024,
available from \url{http://pari.math.u-bordeaux.fr/}.

\bibitem{Shimura1971}
G.~Shimura, \emph{Introduction to the Arithmetic Theory of
Automorphic Functions}, Princeton University Press, 1971.

\bibitem{Silverman1994}
J.~H.~Silverman,
\emph{Advanced Topics in the Arithmetic of Elliptic Curves},
Graduate Texts in Mathematics 151, Springer, 1994.

\bibitem{Silverman2009}
J.~H.~Silverman,
\emph{The Arithmetic of Elliptic Curves},
2nd edition, Graduate Texts in Mathematics 106, Springer, 2009.

\bibitem{TaylorWiles1995}
R.~Taylor, A.~Wiles,
``Ring-theoretic properties of certain Hecke algebras'',
\emph{Ann.\ of Math.} (2) \textbf{141} (1995), no.~3, 553--572.

\bibitem{Wiles1995}
A.~Wiles,
``Modular elliptic curves and Fermat's last theorem'',
\emph{Ann.\ of Math.} (2) \textbf{141} (1995), no.~3, 443--551.

\end{thebibliography}

\end{document}
